\documentclass{situationist-preprint}

\title{Non-Primordial Givenness Suffices}
\msid{PHE-12-a}
\msfamily{PHE-12}
\msdiscipline{Phenomenology}
\mstype{Principal manuscript}
\mskeywords{empathy; Einf\"uhlung; other minds; Husserl; Stein; appresentation; primordial givenness; direct perception; Zahavi; Gallagher; knowledge}

\begin{document}
\maketitle

\begin{abstract}
Does empathy give us other minds? In the phenomenological tradition, empathy (\emph{Einf\"uhlung}) is the experience in which another's mental life is given to me: I see the other's pain in her face, not infer it from her grimace. Husserl and Stein held that this givenness is direct but non-primordial: the other's experience is given as hers, as lived by her, and so not as I live my own. Critics object that a givenness that is not primordial is not really givenness, and that empathy is therefore either inference in disguise or projection. We argue that empathy gives other minds directly though non-primordially, and that non-primordial givenness suffices for knowledge. We distinguish directness, which concerns whether the experience of the other's state is mediated by an inference from something else, from primordiality, which concerns whether the state is given as lived by the one to whom it is given. We show that the second is not required for knowledge, since its absence is not a defect but the form in which another's experience must be given if it is given as another's. Empathy is thus an ordinary source of knowledge of particular minds.
\end{abstract}

\section{The problem}

The traditional problem of other minds asks how I can know that others have experiences, given that I have direct access only to my own. The argument from analogy, in Mill's form, answers that I observe in myself a connection between my experiences and my behavior, observe similar behavior in others, and infer similar experiences \citep{Mill1865}. The theory-theory and simulation theory in cognitive science offer successors: I attribute mental states to others by applying a folk-psychological theory, or by simulating their situation in myself \citep{Goldman2006}. All these views share the assumption that another's mind is never given in experience, only inferred or modeled from what is given.

The phenomenological tradition denies the assumption. Scheler argued that we perceive others' joy in their laughter and their shame in their blushing, and that the argument from analogy misdescribes the phenomena \citep{Scheler1923}. Husserl, in the fifth \emph{Cartesian Meditation}, held that the other's body is given to me as a lived body through an \enquote{appresentation} founded on a pairing with my own, so that the other's experiences are co-given with the body I perceive, as the back of a thing is co-given with its front, though unlike the back they can never be made present to me as they are to the other \citep[\S\S50--54]{Husserl1950}. In \emph{Ideas~II} he describes the other's body as an expressive unity whose mental sense is grasped in its gestures \citep[\S\S43--47]{Husserl1952}. Stein, in her dissertation of 1917, gave the most careful account: empathy is a sui generis act in which another's experience is given to me as non-primordial, that is, as an experience that is primordial for the other but not for me \citep{Stein1917}. \citet{Zahavi2014} and \citet{Gallagher2008} have defended versions of this view in current debates, arguing that we often perceive others' mental states directly.

Two objections dominate. The first says that if the other's experience is given non-primordially, it is not given at all; what is given is the face, and the experience is supplied by me. The second says that even if empathy is a kind of givenness, it cannot ground knowledge, since it is compatible with the other's having no experience at all: the actor's grief, the robot's smile. We argue against both. The question here is distinct from the solipsist's. MND-06 argues that other minds exist, by inference to the best explanation, against one who doubts that there is anything there; the present result is an account of ordinary knowledge of particular minds, which presupposes that the solipsist has been answered or need not be.

\section{Framework}

\begin{definition}[Direct givenness]\label{d:direct}
A state $s$ is given \emph{directly} to a subject in an experience $E$ if $E$ presents $s$ itself, and the subject's taking $s$ to obtain is not the conclusion of an inference from something else presented in $E$.
\end{definition}

\begin{definition}[Primordial givenness]\label{d:prim}
A mental state $s$ is given \emph{primordially} to a subject if it is given as lived through by that subject; it is given \emph{non-primordially} if it is given as lived through by another.
\end{definition}

The distinction is Stein's \citep[ch.~2]{Stein1917}. When I remember a past joy, the joy is given as having been lived by me, though not as now lived; memory is non-primordial with respect to the present but its object is my own. When I empathize with your joy, the joy is given as being lived now, by you. It is present, and it is given to me, but as yours. Stein insists that this is not a deficiency in my grasp but the form of the other's experience as other.

\begin{definition}[Empathy]\label{d:emp}
\emph{Empathy} is an experience in which a mental state of another is given to a subject directly and non-primordially, through the perceived expressive conduct of the other's lived body.
\end{definition}

\begin{principle}[Expressive unity]\label{p:expr}
A person's expressive conduct, such as facial expression, gesture and tone of voice, is not a sign of a mental state separable from it, but in central cases a part or aspect of the state's manifestation: the anger is in the clenched fist and the raised voice.
\end{principle}

Principle~\ref{p:expr} is the thesis Scheler, Husserl and Merleau-Ponty share \citep{Scheler1923,Husserl1952,MerleauPonty1945}. It has an analytic defense: \citet{Green2007} argues that expressions show the states they express, in the sense that perceiving the expression is perceiving part of the state, as seeing the visible part of an object is seeing the object. It does not say that every mental state is expressed, or that every expression is sincere.

\begin{principle}[Knowledge from givenness]\label{p:know}
If a state is given to a subject in an experience that, in normal conditions, is so given only if the state obtains, and the subject believes that the state obtains on the basis of that experience, then, absent defeaters, the subject knows that it obtains.
\end{principle}

Principle~\ref{p:know} is a modest, broadly reliabilist condition on knowledge from experience, of the kind that underwrites perceptual knowledge generally. It is weaker than the claim of EPI-19 that a seeming that $p$ gives prima facie justification for believing that $p$, which it is consistent with.

\section{The argument}

\begin{lemma}[Directness]\label{l:direct}
In central cases, empathy gives another's mental state directly.
\end{lemma}

\begin{proof}
By Principle~\ref{p:expr}, in central cases the perceived expression is a part or aspect of the state's manifestation. So in perceiving it, the subject perceives the state as manifest, not a separate sign from which the state could be inferred. The phenomenology confirms this: I do not see a contraction of facial muscles and infer anger; I see an angry face, and I would have to abstract laboriously to describe the muscles. Nor is there an implicit inference: the inference from analogy requires a premise about my own correlations between experience and expression, which I mostly lack, since I rarely see my own face \citep{Scheler1923}. By Definition~\ref{d:direct}, the state is given directly.
\end{proof}

\begin{lemma}[Non-primordiality]\label{l:nonprim}
Empathy cannot give another's state primordially.
\end{lemma}

\begin{proof}
To be given primordially is to be given as lived by the subject to whom it is given (Definition~\ref{d:prim}). If another's state were given to me as lived by me, it would be given as my state, and so not as another's; the experience would be a case of my having the state, not of empathy. Husserl's version: if what belongs to the other's own essence were directly accessible in the way my own is, it would be merely a moment of my own essence, and the other and I would be the same \citep[\S50]{Husserl1950}. So another's state, given as another's, is given non-primordially.
\end{proof}

\begin{theorem}\label{t:know}
Empathy gives other minds directly though non-primordially, and empathic belief formed in normal conditions, absent defeaters, is knowledge.
\end{theorem}

\begin{proof}
The first conjunct is Lemmas~\ref{l:direct} and~\ref{l:nonprim}. For the second, consider empathy in normal conditions: a person not acting, not concealing, and in a shared form of life in which her expressions have their usual sense. In such conditions an angry face is given as angry only if the person is angry, since by Principle~\ref{p:expr} the expression is part of the anger's manifestation, and manifestation in normal conditions requires that what is manifested obtain. So the experience satisfies the antecedent of Principle~\ref{p:know}, and belief formed on its basis is knowledge absent defeaters. Non-primordiality does not affect this reasoning, since Principle~\ref{p:know} requires that the state be given in a reliably connected experience, not that it be given as lived by the subject. Hence non-primordial givenness suffices for knowledge.
\end{proof}

The theorem's key point is that the condition empathy fails to meet, primordiality, is one no knowledge of another's mind could meet. The demand that the other's experience be given as I live my own is not a higher standard of evidence but a demand that the other be me. A standard that no possible case of knowing another's mind could meet is not a standard for such knowledge. The asymmetry Husserl describes, in which what is not originally accessible is nonetheless accessible in a verifiable way, is the form of knowledge of other minds, not an obstacle to it \citep[\S52]{Husserl1950}.

There is a structural parallel with perception of things. PHE-03 argues that every thing is given in adumbrations whose series never converges on complete givenness: the back of the cube is co-given but never given as the front is. Perceptual knowledge does not require complete givenness. Empathy differs in that the other's experience can never be made present to me even in principle, as the back of the cube can by walking around it, but the principle that knowledge requires reliable givenness, not total givenness, applies to both.

\section{Objections and replies}

\begin{objection}[Projection]\label{o:proj}
What I call perceiving your anger is projecting onto your face a state modeled on my own; the givenness is my construction.
\end{objection}

\begin{reply}
Projection would make empathic experience insensitive to differences between others and myself. But empathy often presents states unlike any I am in, and corrects my expectations: I see that you are hurt by what I thought a joke. Stein emphasizes that empathy can present experiences I have never had and values I do not share, and that it is the basis on which I learn that my own experience is one perspective among others \citep[ch.~3]{Stein1917}. \citet{Zahavi2014} argues that simulation accounts confuse empathy, a way of being directed at another's experience as hers, with emotional contagion or imaginative reproduction, which are about states in me.
\end{reply}

\begin{objection}[Actors and deceivers]\label{o:act}
An actor's grief looks the same as a mourner's. If empathy were knowledge, the audience would know that the actor grieves.
\end{objection}

\begin{reply}
The objection assumes that knowledge from experience requires infallibility. It does not, as perceptual illusions do not refute perceptual knowledge. The actor's performance is a case in which the conditions are not normal, and the audience typically knows this, which defeats any belief that the actor grieves; a skilled deceiver produces a non-normal case that the victim does not recognize, and the victim's belief, though empathically based, is false and so not knowledge. On a disjunctivist view of the kind \citet{McDowell1982} defends for criteria, the deceptive case and the genuine case need not provide the same epistemic standing, since in the genuine case the other's state is itself manifest. Either way, the possibility of error does not show that the good case is not knowledge.
\end{reply}

\begin{objection}[Not perception]\label{o:notperc}
Empathy is not perception, since mental states are not visible properties; calling it givenness is a metaphor.
\end{objection}

\begin{reply}
The thesis does not require that empathy be a sensory perception of a mental property. It requires that the other's state be presented in an experience without inference, which Lemma~\ref{l:direct} establishes. \citet{Smith2010} argues that we see others as minded, in the way we see objects as having hidden sides, and \citet{Gallagher2008} that direct social perception is embedded in interaction rather than detached observation. Whether empathy is called perception is a terminological question; Stein classes it as a sui generis act akin to perception in being direct, unlike it in being non-primordial \citep{Stein1917}. The result requires only the directness.
\end{reply}

\begin{objection}[Hidden states]\label{o:hidden}
Many of another's states are not expressed: her thoughts while she sits still, her concealed resentment. Empathy cannot reach them, so it cannot be our way of knowing other minds.
\end{objection}

\begin{reply}
The thesis is that empathy gives other minds, not that it gives all of another's mind. Unexpressed states are known, when they are, by testimony, by inference from conduct over time, or not at all. Empathy is the basic and ordinary form of access, the one presupposed when we interpret what others say and do; it need not be the only one.
\end{reply}

\section{Conclusion}

Empathy gives other minds directly, since the other's state is manifest in expressive conduct and is not inferred from it, and non-primordially, since it is given as lived by the other and not by me. Non-primordial givenness suffices for knowledge. The condition it fails to meet, primordiality, is one no knowledge of another's mind could meet without the other ceasing to be other. What knowledge requires is that the other's state be given in an experience reliably connected to it, and empathy in normal conditions satisfies that requirement.

\begin{thebibliography}{99}\small
\bibitem[Gallagher(2008)]{Gallagher2008} Gallagher, S. (2008). Direct perception in the intersubjective context. \emph{Consciousness and Cognition}, 17(2), 535--543.
\bibitem[Goldman(2006)]{Goldman2006} Goldman, A.~I. (2006). \emph{Simulating Minds: The Philosophy, Psychology, and Neuroscience of Mindreading}. Oxford University Press.
\bibitem[Green(2007)]{Green2007} Green, M.~S. (2007). \emph{Self-Expression}. Oxford University Press.
\bibitem[Husserl(1950)]{Husserl1950} Husserl, E. (1950). \emph{Cartesianische Meditationen und Pariser Vortr\"age}. Husserliana I. Ed. S. Strasser. Martinus Nijhoff. Trans. D. Cairns, \emph{Cartesian Meditations}, Martinus Nijhoff, 1960.
\bibitem[Husserl(1952)]{Husserl1952} Husserl, E. (1952). \emph{Ideen zu einer reinen Ph\"anomenologie und ph\"anomenologischen Philosophie. Zweites Buch}. Husserliana IV. Ed. M. Biemel. Martinus Nijhoff. Trans. R. Rojcewicz and A. Schuwer, \emph{Ideas Pertaining to a Pure Phenomenology and to a Phenomenological Philosophy, Second Book}, Kluwer, 1989.
\bibitem[McDowell(1982)]{McDowell1982} McDowell, J. (1982). Criteria, defeasibility, and knowledge. \emph{Proceedings of the British Academy}, 68, 455--479.
\bibitem[Merleau-Ponty(1945)]{MerleauPonty1945} Merleau-Ponty, M. (1945). \emph{Ph\'enom\'enologie de la perception}. Gallimard. Trans. D.~A. Landes, \emph{Phenomenology of Perception}, Routledge, 2012.
\bibitem[Mill(1865)]{Mill1865} Mill, J.~S. (1865). \emph{An Examination of Sir William Hamilton's Philosophy}. Longmans, Green.
\bibitem[Scheler(1923)]{Scheler1923} Scheler, M. (1923). \emph{Wesen und Formen der Sympathie}. Friedrich Cohen. Trans. P. Heath, \emph{The Nature of Sympathy}, Routledge \& Kegan Paul, 1954.
\bibitem[Smith(2010)]{Smith2010} Smith, J. (2010). Seeing other people. \emph{Philosophy and Phenomenological Research}, 81(3), 731--748.
\bibitem[Stein(1917)]{Stein1917} Stein, E. (1917). \emph{Zum Problem der Einf\"uhlung}. Buchdruckerei des Waisenhauses. Trans. W. Stein, \emph{On the Problem of Empathy}, ICS Publications, 1989.
\bibitem[Zahavi(2014)]{Zahavi2014} Zahavi, D. (2014). \emph{Self and Other: Exploring Subjectivity, Empathy, and Shame}. Oxford University Press.
\end{thebibliography}

\end{document}
